\couple{Sensei}{ある学生が禅師に自己紹介した。
「写真家はどうやって自分の写真を撮るのですか？

学生は笑った。
先生は棒で生徒を叩いた。
生徒は呻いた。
先生は棒で生徒を殴った。
生徒は泣いた。
先生が棒で生徒を殴った。

生徒の目が輝いた
彼は餓死するまで動かずに立っていた}{A student introduced himself to the Zen Master.
“How does a photographer take a picture of himself?” the teacher asked.

The student laughed.
The teacher hit the student with a stick.
The student grunted.
The teacher hit the student with a stick.
The student cried.
The teacher hit the student with a stick.

The student's eyes light up.
He stood motionless until he starved to death}

\couple{Iara}{Proyecciones no invertibles

Pensaba que veía a los otros
Pensaba que me veía al mirar a los otros

Ahora intento verme mirando a los otros que se miran a sí mismos mirándome
Pero veo menos

Si continúo
Ya no veo

Soy otro
Que no ve
Pero mira

Me mira y se mira a sí mismo}{Non-Invertible Projections

I thought I was seeing others
I thought I was seeing myself by looking at others

Now I try to see myself by looking at the others who are looking at themselves by looking at me
But I see less

If I continue
I no longer see

I am other 
Who does not see
But looks

Looks at me and looks at himself}

\couple{Stili Crudi}{Sorrido al fotografo

Senile, il canto dell’alba
Orchestra il tempo
Rimasto
Rimpianto
In teoria mai perduto ma
D’altro canto inaccessibile.
Odo versi che depongono 
Altri significati
Lodi e lamenti al giallo
Fiore dell’avvenire.
Onde di luce
Tormentano
Ogni 
Glifo 
Rampante
Allocato a costruire 
Fantastiche statue di sale
Onta alla dolce complessità umana}{I smile at the photographer

Senile, the song of dawn
Orchestrates time
Remaining
Regret
In theory never lost but
On the other hand inaccessible.
I hear verses that lay down
Other meanings
Praises and laments to the yellow
Flower of the future.
Waves of light
Torment
Each
Glyph
Rampant
Allocated to build
Fantastic statues of salt
Shame on the sweet human complexity}

\couple{W. Worseworth}{Mi stavo appropinquiando al balcone, dopo una giornata di lavoro, con in mente una bella storia distopica sull'invenzione di un dispositivo in grado di fare foto multidimensionali: capace di imprimere tutte le (-infinite-) finite (-possibili-) attuali sfumature di una persona; e di come, ben presto integrato in tutti i cellulari, avesse poi cambiato il mondo; dal fondo delle masse fino in cima ai potenti. 
Ma come ho messo il naso fuori dalla porta-finestra mi ha colpito subito una piacevolissima aria primaverile. Dopo una settimana di piogge il cielo aveva riconquistato i colori tipici dell'imbrunire, con le ultime tracce di luce che risaltavano tratti delle colline in lontananza. I primi pipistrelli, che sono poi il corrispettivo notturno delle rondini, mi svolazzavano davanti e gli alberi di fronte casa passavano, da sinistra verso destra, dal verde acceso delle nuove, ritrovate gemme al secco dell'inverno a seconda della diversa esposizione al sole. I bar sui quali il balcone affacciava avevano iniziato a lasciare i tavolini fuori fino a tardi e un vociare di persone faceva da sottofondo al canto di qualcheduno di quegli uccelli con la coda allungata di cui qua è pieno. 

Mentre una folata di profumo proveniente dalla cucina mi ricordava la pizza a scaldare nel forno ho pensato che la realtà che ci circonda, nella sua piu brutale semplicità, sta lì come un salvagente in mezzo al mare. Ci basta afferrarlo un attimo, qualche secondo, per riconnetterci con noi stessi e prendere fiato dall'immondizia che la nostra mente ci propina ogni giorno fingendosi giudice, nemico o paparazzo.}{I was approaching the balcony after a long day of work, with a nice dystopian story in mind about the invention of a device capable of taking multidimensional photos: able to capture all the (-infinite-) finite (-possible-) actual shades of a person; and how, once quickly integrated into all mobile phones, it had soon changed the world; from the depths of the masses to the heights of the world's powerful.

But as soon as I stepped outside the glass door, I was immediately struck by a delightful spring breeze. After a week of rain, the sky had regained the typical colors of dusk, with the last traces of light highlighting the contours of the distant hills. The first bats—who are, after all, the nighttime counterpart of swallows—fluttered before me, while the trees in front of my house shifted from the bright green of newly sprouted buds to the dryness of winter, depending on their exposure to the sun. The bars overlooking my balcony had begun leaving their tables outside until late, and the hum of conversations blended with the singing of some of those birds with long tails that are so common around here.

As a wave of scent from the kitchen reminded me of the pizza warming in the oven, I thought about how the reality surrounding us, in its most brutal simplicity, stands there like a lifebuoy in the middle of the sea. All we need is to grasp it for a moment, just a few seconds, to reconnect with ourselves and catch our breath from the garbage our mind feeds us every day, disguising itself as a judge, an enemy, or a paparazzo.}

\couple{Il Barone Tartaglia }{Ho visto il viso, mio, ripetuto
circondarmi la vista. 

Molti appesi ai muri, distesi in sequenza, 
guardare la strada attraverso schermi bianchi. 
Altri su palloncini colorati, volanti,
frantumati nei vetri dei piani alti
guardare l'orizzonte scivolare come luce sul fiume. 
Pochi su di enormi cartelloni di legno fradicio,
innalzandosi oltre mari filiformi di strade e utilitarie,
guardare con occhi liquidi l'inaccessibile moltiplicarsi.
Alcuni su televisioni esposte,
apparire e scomparire in loop elettromagnetici e sciabordare di elettroni, 
visi guardati e non guardanti, 
telematica assenza di realta'. 
I visi tutti ripetono la stessa espressione 
come forme iperplastiche, mutevoli e immobili, 
di vetro. 

Circondato dicevo, 
non posso evitarne gli sguardi 
che si allineano col mio in se stessi,
come fasci di luce di fari distanti
che per un istante si congiungono 
nell'illusione dello specchio. 
La fine coincide con l'inizio, eppure esiste una distanza,
una menzognera alterazione dello spazio.
I visi penetrano dentro di me,
bacio e stupro degli occhi, 
divorandomi e trascinandomi fuori come uno sputo, 
un eco di rimando, risonante, che fugge da perfette coincidenze di forma. 
Sempre vivido nella memoria 
il tuonante suono del flash, quando la luce mi investiva 
e lo spettro del mio viso veniva catturato dall'obiettivo. 
E oltre i confini del bianco, ancora una voce crudele dice
"Sorridi, presto tutti ti vedranno".

Sento la mia vita strapparsi come la carne del pollo,
che morsa tra denti abbandona l'osso, 
unico avanzo del corpo gia' perduto.


 


}{I saw my face, my own, repeated,
surrounding my sight.

Many, hanging on walls, stretched in sequence,
watching the street through white screens.
Others, on floating, colored balloons,
shattered in the glass of high floors,
watching the horizon slip away like light on the river.
Few, on vast, rotting wooden billboards,
rising beyond threadlike seas of roads and small cars,
gazing with liquid eyes at the inaccessible multiplying.
Some, displayed on televisions,
appearing and vanishing in electromagnetic loops,
in the sloshing of electrons,
faces seen but never seeing,
a telematic absence of reality.

All faces repeat the same expression,
like hyperplastic forms, mutable yet still,
of glass.

Surrounded, I said,
I cannot escape their gazes,
which align with mine within themselves,
like beams of distant headlights
that for an instant merge
in the illusion of a mirror.

The end coincides with the beginning, yet there is a distance,
a deceitful distortion of space.
The faces penetrate me,
a kiss and a rape of the eyes,
devouring me, dragging me outward like spit,
an echo rebounding, resonating,
fleeing from perfect coincidences of form.

Always vivid in memory,
the thundering sound of the flash, when the light struck me
and the specter of my face was captured by the lens.
And beyond the borders of the white,
a cruel voice still whispers,
"Smile, soon everyone will see you."

I feel my life tearing apart,
like the flesh of a chicken,
which, bitten by teeth, abandons the bone,
the only remnant of a body already lost.}

\couple{Tom Joad}{qualunque cristallizzazione di un istante sarà sempre capace di immortalare la vera natura del suo soggetto in quell'istante nel contesto in cui viene operata.
la maschera rimarrà una maschera, non muterà mai in quanto tale, in quanto volontaria assunzione di un'espressione finta e costruita.
il giullare farà ridere, in quanto indipendentemente dall' espressione da egli assunta, porterà con sé la sua placida inettitudine.
colui che è triste non potrà camuffare il lamento del suo cuore con alcun artificio, 
colui che è felice non potrà che apparire come degno protagonista o come sfrontato egocentrico.
il conformista sparirà tramutandosi in oggetto di scena, troppo educato per opporsi al comando dell architetto,
l anarchico rovinerà l opera agli occhi del conformista,
il furbo distruggerà il sogno dell anarchico di stupire l architetto,
il malvagio farà la sua solita faccia da culo.
prepararsi per quel momento ha la stessa utilità di sapere prima di avere una malattia incurabile, il processo rimane conservativo.
chi vive per gli altri dedicherà agli altri la sua posa e il suo sorriso,
c'è  chi vive nella foto per essere ammirato, così come fa sempre. c'è anche chi vive perché la vita esiste,
la loro posa sarà imprevedibile in quanto puramente contestuale.
se sorrido è per non fare un torto all'architetto, ma non cambierò il mondo di quell istante.
domani forse non ci sarà più nulla e rimarrà solo quella foto.
allora meglio che mi giri e cali le braghe}{Any crystallization of a moment shall always be capable of immortalizing the true nature of its subject at that very instant, within the context in which it is carried out. The mask shall remain a mask, never changing in its essence, for it is a voluntary adoption of a false and constructed expression. The jester will make one laugh, for regardless of the expression he assumes, he will carry with him his placid ineptitude. He who is sorrowful will not be able to disguise the lament of his heart with any artifice. He who is happy can only appear either as a worthy protagonist or as a brazen egocentric. The conformist shall vanish, transforming into a mere object in the scene, too well-mannered to oppose the architect’s command. The anarchist will ruin the work in the eyes of the conformist. The cunning one will destroy the anarchist’s dream of astonishing the architect. The wicked will wear his usual contemptuous face.

To prepare for such a moment serves the same purpose as knowing beforehand that one has an incurable disease; the process remains preservative. He who lives for others will dedicate his pose and his smile to them. There are those who live for the photo to be admired, just as they always do. There are also those who live because life exists. Their pose will be unpredictable, for it is purely contextual. If I smile, it is to avoid doing the architect an injustice, but I will not change the world of that moment. Tomorrow, perhaps, there will be nothing left, and only that photo will remain. Then, it is better I turn around and lower my pants.



}

\couple{Pirlandello}{Cosa vuole da me questa figura?
Mi guarda negli occhi, non si muove.
Avrà trovato un altro pensiero da leggere in silenzio.
Nel mio sguardo.
Ennesimo soffice abbraccio del nulla.

Cosa vuole da me questa figura?
Mi guarda negli occhi, devo fare lo stesso.
Come sono caduto in questo tacito gioco di sguardi?
Forse avevo altre cose da fare.
Forse no.

Cosa vuole da me questa figura?
E se non esistesse? Provo a toccarla.
Come io mi muovo, lei già si era mossa.
Congiunte le punte dei nostri duplici indici. Nulla accade.
Avevo sicuramente altre cose da fare.

Cosa vuole da me questa figura?
Se lei è, io sono. Lei è, se io sono.
Io sono un pugno.
E lei mille pugni, sporchi del mio sangue, sporchi del suo sangue.
Un sangue vitreo che riflette mille copie della stessa identica faccia.

Cosa vuole da me questa figura?
Una folla di occhi impressi nel suo disgregarsi.
Un unicità che si spegne in una molteplicità di sguardi.
Serve troppo sangue per frammentarne i frammenti.
Mi resta il sangue ed uno specchio rotto.

Una goccia di sangue cade sul foglio.
La pace parcheggia in divieto di sosta. 
Tra me e i frammenti di lei.
I suoi mille sguardi diventan sorrisi.
Ed io sorrido.}{What does this figure want from me?
It stares into my eyes, unmoving.
Has it found another thought to read in silence?
In my gaze.
Yet another hollow embrace of nothing.

What does this figure want from me?
It stares into my eyes, I must stare back.
How did I fall into this silent game of glances?
Perhaps I had other things to do.
Perhaps not.

What does this figure want from me?
What if it isn’t real? I reach to touch it.
As I move, it had already moved.
Our twin index fingers meet. Nothing happens.
I definitely had other things to do.

What does this figure want from me?
If it is, then I am. It is, if I am.
I am a fist.
And it, a thousand fists, smeared with my blood, smeared with its blood.
A glassy blood that reflects a thousand copies of the same face.

What does this figure want from me?
A crowd of eyes etched into its dissolution.
A singularity drowned in a swarm of stares.
Too much blood is needed to shatter the shards.
All I’m left with is blood and a broken mirror.

A drop of blood falls onto the page.
Peace parks in a no-standing zone.
Between me and the fragments of it.
Its thousand stares twist into smiles.
And I smile back.}

\couple{Una mucca con il corpo di Saturno e i piedi da umano}{Un ulteriore esempio di sciatto lavoro di critica si puo trovare nelle pubblicazione
di Caca Peerley. Come spesso accade nelle biografie di critici scadenti, Peerley
inizio la sua carriera letteraria cercando di pubblicare i propri racconti ma senza
mai riuscire ad trovare un proprio pubblico o, per quel che conta, un proprio
stile. Entrato in accademia e quindi in ambienti di critica da salotto riuscı a
trovare la propria vocazione: pubblicare editoriali di critica letteraria con
l’apparente ambizione di riuscire a criticare sempre nel modo piu insignificante ed
invidioso. Di seguito e riportato un estratto esemplare da ”Poetica dell’inetto”
(Mortemi Editore, Roma, 1998): 

—-Le pochissime opere di Pirka Kalei hanno raggiunto una notorieta dal
sottoscritto incomprensibile data la loro mancanza di ogni tipo di significato e
la sua ultima impresa pubblicata postuma ed incompiuta non sembra da meno.
Questo scritto mai dichiarato pronto alla pubblicazione rappresenta perfetta-
mente la pusillaminita dell’autore. ”” 
[Folla:] Ciao, ci piace camminare per i
corridoi e sorridere.
[Cane:] A me piace fare le cose a caso. Forse perche sono un cane. Forse
perche non sono un cane.
[Folla:] Ma che bel libro di inglese! Ahah
[Cane:] Sia io che voi sappiamo che questo non centrava nulla.
[Folla:] Gli ultimi corti di Wes Anderson sono un vero capolavoro. Riescono
ad essere piu Wes Anderson di Wes Anderson. Ve li consiglio.
[Cane:] Si li ho visti ma non me li ricordo molto. Pero sı hai ragione sono
molto belli.
[Folla:] Avete visto cosa ha pubblicato Trump? Non so come fanno gli amer-
icani ad essere cosı stupidi. No ma e l’Europa che dovrebbe riarmarsi.
[Cane:] Ma mi lasciate stare per favore? Stavo cercando di scrivere sta roba
e gia faccio fatica e tutte ste voci mi incasinano la testa.
[Folla:] Vuoi venire al mare domani? Ci dovrebbe essere il sole.
[Cane:] Eh si avrei da fare e da finire questa cosa. Non so se riesco a finire
oggi. Ma si. Massı che ci vengo dai. Certo che ci voglio venire. Non e un
problema. Ci vediamo lı eh.
[Folla:] Guarda un flash!
[Cane:] Occristo
[Cani:] In questa spiaggia non ci sono bidoni. Ma guarda che la raccolta
differenziata e una truffa. Viene riciclato solo il 5percento della plastica ed il
resto finisce nelle discariche in Asia o direttamente nel mare. Lo sapevate che
Elly’s chiude? Il bar con con i tavoli fatti di yoghurt? Ci ero andato solo per
un date con una tipa. Peccato sembrava bello. Riciclare non e la soluzione alla
sovrapproduzione di plastica.
[Cane:] Comunque sono proprio belli i film di Wes Anderson eh. Tutte
quelle cose alla Wes Anderson e poi gli ultimi corti eh che roba.
[Cani:] Guarda un flash!
[Cane:] Occazzo
[Cane:] Ma la volete smettere di fare tutto sto casino? Sto cercando di fare
qualcosa di utile qui! Che una pagina e gia troppo poco!
[Cane:] Ma con chi parli?
[Cane:] Con loro! E con te! Come fate a scodinzolare tra tutti questi flash?
Tra tutte queste opinioni? Ridere leggere o parlare senza mai un dubbio? Me
ne vado.
[Cane:] Cane no dai, tranquillo, sorridi. ””
Leggendo e evidente la motivazione per cui Kalei non si decise mai a pubbli-
care questo racconto il cui cui titolo sarebbe dovuto essere ’Sorrido al fotografo’.
Dal testo stesso traspare la meschinita che pervade l’interezza della sua opera.
—--

Un’ analisi cosı miope e malposta puo venire solo da un’ anima amara come
quella di Peerley e mi auguro che i futuri critici che stanno leggendo questa
antologia riescano a mantenere intatta la propria relazione con la realta meglio
di quanto non ci si riuscito questo rinomato commentatore.
[Folla:] Ma che problemi aveva quel cane?
[Io:] Non so ma quello che ho scritto mi fa cagare.}{Another example of sloppy critical work can be found in the publications of Caca Peerley. As often happens in the biographies of mediocre critics, Peerley began his literary career trying to publish his own short stories, but never managed to find an audience—or, for that matter, a style. After entering academia and the world of drawing-room criticism, he discovered his true calling: publishing literary editorials with the apparent ambition of criticizing in the most petty and envious way possible. Below is a representative excerpt from Poetics of the Inept (Mortemi Publishing, Rome, 1998):

“The very few works of Pirka Kalei have reached a level of notoriety that I personally find incomprehensible, given their utter lack of meaning in any form. His last project, published posthumously and unfinished, is no different. This piece—never declared ready for publication—perfectly embodies the author’s cowardice.

[Crowd:] Hi, we like walking through the halls and smiling.
[Dog:] I like doing random things. Maybe because I’m a dog. Maybe because I’m not a dog.
[Crowd:] What a lovely English textbook! Haha
[Dog:] You and I both know this has nothing to do with anything.
[Crowd:] Wes Anderson’s latest shorts are masterpieces. They manage to be more Wes Anderson than Wes Anderson himself. Highly recommend.
[Dog:] Yeah, I saw them but don’t remember much. Still, you’re right—they’re really good.
[Crowd:] Have you seen what Trump posted? I don’t get how Americans can be so stupid. But no, it’s Europe that needs to rearm.
[Dog:] Can you all please leave me alone? I’m trying to write this thing, and I’m already struggling—these voices are messing with my head.
[Crowd:] Wanna go to the beach tomorrow? Should be sunny.
[Dog:] Ugh, I’ve got stuff to do… need to finish this… Don’t know if I’ll manage today. But yeah. What the hell, I’ll come. Of course I want to come. No problem. See you there.
[Crowd:] Look, a flash!
[Dog:] Oh Christ.
[Dogs:] There are no bins on this beach. But you know recycling is a scam, right? Only 5\% of plastic gets recycled, the rest ends up in landfills in Asia or straight into the ocean. Did you hear Elly’s is closing? The bar with yogurt tables? I only went there once on a date. Shame, seemed nice. Recycling won’t fix plastic overproduction.
[Dog:] Anyway, Wes Anderson’s movies are really beautiful, huh. All that Wes Anderson stuff, and those latest shorts—what a trip.
[Dogs:] Look, a flash!
[Dog:] Oh shit!
[Dog:] Can you all stop making such a racket? I’m trying to do something useful here! A single page is already too little!
[Dog:] Who are you talking to?
[Dog:] To them! And to you! How can you wag your tails through all these flashes? Through all these opinions? Laughing, reading, talking—never a single doubt? I’m out.
[Dog:] Don’t go, Dog. Chill out. Smile.”

Reading it, the reason why Kalei never chose to publish this story—titled Smiling at the Photographer—becomes painfully clear. The text reeks of the same pettiness that pervades his entire body of work.

Such a short-sighted and misplaced analysis could only come from a bitter soul like Peerley’s, and I sincerely hope that the future critics reading this anthology manage to preserve their relationship with reality better than this celebrated commentator ever could.

[Crowd:] What was wrong with that dog?
[Me:] I don’t know, but what I wrote sucks.}

